\section{Introduction}
\label{sec:introduction}

Nuclear fission is intrinsically a quantum-mechanical process. The formation of an excited compound nucleus, the evolution of nuclear degrees of freedom toward scission, quantum tunnelling through the fission barrier, and the subsequent formation of fragments are all governed, at the microscopic level, by quantum dynamics. Modern microscopic approaches therefore describe the nucleus as a complex interacting many-body quantum system and employ nuclear density functional theory, time-dependent methods, and collective quantum dynamics to investigate the path from excitation to scission \cite{SchunckRobledo2016}.

At the reactor-engineering scale, however, this microscopic quantum dynamics is not treated as an active control degree of freedom. Reactor physics is formulated primarily in terms of neutron transport, reaction rates, isotopic concentrations, thermal-hydraulic variables, and macroscopic feedback mechanisms. The underlying nuclear dynamics enters through quantities such as reaction cross sections, decay constants, neutron multiplicities, and evaluated nuclear data. Consequently, the internal quantum state of the fissile nuclei is normally not regarded as a dynamically controllable reactor variable.

Reactor control therefore acts predominantly on the environment in which fission takes place rather than on the nuclear state from which fission proceeds. Control rods modify neutron absorption; moderators alter neutron energy distributions; soluble absorbers modify macroscopic absorption; coolant and fuel temperatures provide reactivity feedback; and changes in geometry or composition affect neutron leakage and multiplication. In each case, the principal strategy is to control the neutron population or its interaction with an effectively passive fissile medium.

This situation differs from that encountered in atomic and optical physics. Coherent control makes it possible not only to predict transition probabilities but also to manipulate quantum states by controlling the amplitude, frequency, timing, and phase of applied electromagnetic fields.

Related approaches to the representation and manipulation of controlled
quantum states have also been discussed in the context of quantum-system
design and simulation \cite{SisiniPC}.

Nuclear quantum optics has progressively extended related concepts to nuclear degrees of freedom. Early theoretical work explicitly identified preparation, control, and detection of nuclear states as a possible objective of resonant laser--nucleus interactions \cite{Buervenich2006}.

Subsequent experiments have demonstrated that this idea is physically meaningful in selected systems. Coherent X-ray-optical manipulation of collective nuclear excitations has been demonstrated in $^{57}$Fe \cite{Heeg2021}. Spectroscopic access to the exceptionally low-energy isomeric transition in $^{229}$Th has further shown that nuclear transitions can be addressed with optical-frequency techniques under favorable conditions \cite{Yamaguchi2024}. At much higher photon energies, theoretical work has investigated coherent MeV photon fields interacting with heavy nuclei in regimes where excitation competes with internal equilibration, neutron evaporation, and fission \cite{Palffy2015}.

These developments motivate a question that is distinct from both conventional reactor physics and conventional nuclear spectroscopy:

\begin{quote}
Can externally controlled preparation or manipulation of nuclear states in fissile matter modify the subsequent probability or dynamics of neutron-induced fission in a way that can be represented as a controllable physical variable?
\end{quote}

The question is not whether electromagnetic radiation can itself induce fission, nor whether fission is fundamentally quantum mechanical. Rather, the issue considered here is whether the \emph{state dependence} of the fission process can be exploited as a control mechanism.

Conventionally, a fission rate may be represented schematically as
\begin{equation}
R_f
=
N_f
\int
\sigma_f(E)\,
\Phi(E)\,
dE,
\label{eq:conventional_fission_rate_intro}
\end{equation}
where $N_f$ is the number density or number of fissile nuclei appropriate to the chosen normalization, $\Phi(E)$ is the neutron flux, and $\sigma_f(E)$ is the neutron-induced microscopic fission cross section.

Within this description, reactor control primarily modifies the neutron field, while $\sigma_f(E)$ is treated as a material response supplied by nuclear data. The conceptual extension investigated in this work is to allow the effective fission response to depend explicitly on a nuclear-state distribution,
\begin{equation}
\sigma_f
\rightarrow
\sigma_f(E,\rho_N),
\label{eq:sigma_state_intro}
\end{equation}
where $\rho_N$ denotes a nuclear density operator or, in the incoherent limit, a statistical distribution over nuclear states.

If $\rho_N$ can be influenced by an external control variable $u_Q(t)$,
\begin{equation}
u_Q(t)
\longrightarrow
\rho_N(t)
\longrightarrow
\sigma_f^{\mathrm{eff}}(t)
\longrightarrow
R_f(t),
\label{eq:central_chain_intro}
\end{equation}
then the microscopic state of the fissile material becomes, in principle, an additional degree of freedom in a controlled fission system.

This possibility should not be interpreted as a proposal for an immediately realizable ``quantum nuclear reactor.'' Significant obstacles are evident, including state-preparation efficiency, competing decay channels, rapid equilibration of highly excited nuclei, decoherence, the fraction of nuclei that can be addressed, and the energetic cost of actuation. The purpose of this work is more limited and more fundamental: to formulate the problem as a consistent dynamical control framework and to identify the minimum physical conditions required for nuclear-state manipulation to influence a macroscopic fission system.

The central hypothesis may therefore be stated as
\begin{equation}
\boxed{
\begin{minipage}{0.86\columnwidth}
\centering
The nuclear state of fissile matter, normally implicit in reaction data,\\
may in principle become an explicit control variable.
\end{minipage}
}
\end{equation}

The remainder of this paper develops this statement from conventional reactor control to a minimal microscopic model, identifies experimentally established state-dependent fission in $^{235m}$U as a concrete anchor, and derives a falsifiable path from nuclear-state actuation to reactor-level observables.
